\section{Vectorgeometrie}
Ein Vektor \(\vec{v}\) beschreibt eine Translation (Verschiebung) im Raum und wird durch seine Richtung und Länge \(|\vec{v}|\) charakterisiert.\\
Ein Vektor von Punkt \(A\) zu Punkt \(B\) wird als \(\overrightarrow{AB}\) bezeichnet. Spezielle Vektoren:
\begin{itemize}
     \item \textbf{Nullvektor} \(\vec{0}\): Ein Vektor mit Länge 0, der keine Richtung hat.
     \item \textbf{Einheitsvektor} \(\vec{e}_i\): Ein Vektor, der in der Richtung der \(i\)-ten Achse liegt und die Länge 1 hat.
     \item \textbf{Gegenvektor} \(-\vec{v}\): Ein Vektor, der die gleiche Länge wie \(\vec{v}\) hat, aber in die entgegengesetzte Richtung zeigt.
\end{itemize}

\subsection{Vektoraddition}
Werden zwei Verschiebungen hintereinander ausgeführt, entspricht dies der Addition der Vektoren.

Rechenregeln:
\begin{align*}
     \vec a + \vec 0            & = \vec a                                                     \\
     \vec a + (-\vec a)         & = \vec 0                                                     \\
     \vec a + \vec b            & = \vec b + \vec a \quad \text{(Kommutativgesetz)}            \\
     (\vec a + \vec b) + \vec c & = \vec a + (\vec b + \vec c) \quad \text{(Assoziativgesetz)}
\end{align*}

\subsection{Skalare Multiplikation}
Die Multiplikation eines Vektors mit einer reellen Zahl $\lambda$ verändert seine Länge:

\begin{itemize}
     \item $|\lambda \vec a| = |\lambda| \cdot |\vec a|$
     \item $\lambda > 0$: gleiche Richtung
     \item $\lambda < 0$: entgegengesetzte Richtung
     \item $\lambda = 0$: Nullvektor
\end{itemize}

Wichtige Rechenregeln:
\begin{align*}
     1\vec a                  & = \vec a                        \\
     \lambda(\mu \vec a)      & = (\lambda\mu)\vec a            \\
     (\lambda + \mu)\vec a    & = \lambda\vec a + \mu\vec a     \\
     \lambda(\vec a + \vec b) & = \lambda\vec a + \lambda\vec b
\end{align*}

\subsection{Linearkombination}
Eine Linearkombination von Vektoren $\vec a_1,\dots,\vec a_n$ ist ein Ausdruck der Form

\[
     \lambda_1 \vec a_1 + \lambda_2 \vec a_2 + \dots + \lambda_n \vec a_n \quad \text{mit} \quad \lambda_i \in \mathbb{R}
\]



\subsection{Kollinearität und Komplanarität}

\textbf{Kollinear:}
Zwei Vektoren sind kollinear, wenn einer ein Vielfaches des anderen ist:

\[
     \vec a = \lambda \vec b
\]

\textbf{Komplanar:}
Drei Vektoren sind komplanar, wenn sie in derselben Ebene liegen.
Ist $\vec a$ und $\vec b$ nicht kollinear, kann jeder dritte Vektor dargestellt werden als:

\[
     \vec c = \lambda \vec a + \mu \vec b
\]

Sind drei Vektoren im $\mathbb{R}^3$ nicht komplanar, kann jeder Vektor geschrieben werden als

\[
     \vec d = \lambda \vec a + \mu \vec b + \nu \vec c
\]

\subsection{Koordinatendarstellung}

Mit Einheitsvektoren $\vec e_1, \vec e_2, \dots, \vec e_n$ in $\mathbb{R}^n$ gilt:

\[
     \vec a = a_1 \vec e_1 + a_2 \vec e_2 + \dots + a_n \vec e_n = \sum_{i=1}^{n} a_i \vec e_i
\]

Komponentenschreibweise:

\[
     \vec a =
     \begin{pmatrix}
          a_1    \\
          a_2    \\
          \vdots \\
          a_n
     \end{pmatrix}
\]

Ein Punkt $P(x_1, x_2, \dots, x_n)$ hat den Ortsvektor

\[
     \vec r(P) =
     \begin{pmatrix}
          x_1    \\
          x_2    \\
          \vdots \\
          x_n
     \end{pmatrix}
\]

\subsection{Betrag eines Vektors}

In $\mathbb{R}^n$:

\[
     |\vec a| = \sqrt{a_1^2 + a_2^2 + \dots + a_n^2} = \sqrt{\sum_{i=1}^{n} a_i^2}
\]

\subsection{Vektor zwischen zwei Punkten}

Für Punkte $P(x_{P,1}, x_{P,2}, \dots, x_{P,n})$ und $Q(x_{Q,1}, x_{Q,2}, \dots, x_{Q,n})$ in $\mathbb{R}^n$ gilt:

\[
     \overrightarrow{PQ} = r(Q) - r(P) =
     \begin{pmatrix}
          x_{Q,1} - x_{P,1} \\
          x_{Q,2} - x_{P,2} \\
          \vdots            \\
          x_{Q,n} - x_{P,n}
     \end{pmatrix}
\]

Der Abstand der Punkte ist:

\[
     |PQ| = \sqrt{\sum_{i=1}^{n} (x_{Q,i}-x_{P,i})^2}
\]



\subsection{Skalarprodukt}

Das Skalarprodukt zweier Vektoren $\vec a$ und $\vec b$ ist
\[
     \vec a \cdot \vec b = |\vec a|\,|\vec b|\cos(\varphi),
\]
mit $0 \le \varphi \le \pi$. Das Ergebnis ist eine reelle Zahl.

\subsubsection{Eigenschaften}
Für $\vec a,\vec b,\vec c \in \mathbb{R}^n$ und $\lambda \in \mathbb{R}$:

\begin{itemize}
     \item $\vec a \cdot \vec b = 0 \iff \vec a$ und $\vec b$ sind orthogonal.
     \item $\vec a \cdot \vec a = |\vec a|^2$
     \item $\vec a \cdot \vec b = \vec b \cdot \vec a$
     \item Distributiv in beiden Argumenten:
           \[
                \vec a \cdot (\vec b + \vec c) = \vec a \cdot \vec b + \vec a \cdot \vec c
           \]
     \item Verträglich mit Skalaren:
           \[
                \lambda (\vec a \cdot \vec b) = (\lambda \vec a) \cdot \vec b = \vec a \cdot (\lambda \vec b)
           \]
\end{itemize}

\subsubsection{Berechnung über Komponenten}

In $\mathbb{R}^2$:
\[
     (a_1,a_2)\cdot(b_1,b_2)=a_1b_1+a_2b_2
\]

In $\mathbb{R}^3$:
\[
     (a_1,a_2,a_3)\cdot(b_1,b_2,b_3)=a_1b_1+a_2b_2+a_3b_3
\]

\subsubsection{Winkel zwischen zwei Vektoren}
% TODO: Grafik einfügen: Zwei Vektoren mit Zwischenwinkel phi und Projektion auf die Achse.

\[
     \varphi = \arccos\left(\frac{\vec a \cdot \vec b}{|\vec a|\,|\vec b|}\right)
\]

\subsubsection{Orthogonale Projektion}

Die Projektion von $\vec b$ auf $\vec a$:

\[
     \vec b_a = \frac{\vec a \cdot \vec b}{|\vec a|^2}\,\vec a
\]

und ihre Länge:

\[
     |\vec b_a|=\frac{|\vec a \cdot \vec b|}{|\vec a|}
\]

% TODO: Grafik einfügen: Zerlegung von b in parallelen und orthogonalen Anteil zu a.

\subsection{Vektorprodukt}

Das Vektorprodukt im $\mathbb{R}^3$ liefert einen Vektor $\vec a \times \vec b$.

\begin{itemize}
     \item Betrag:
           \[
                |\vec a \times \vec b| = |\vec a|\,|\vec b|\sin(\varphi)
           \]
     \item $\vec a \times \vec b$ ist orthogonal zu $\vec a$ und $\vec b$
     \item Richtung nach der Rechte-Hand-Regel
\end{itemize}

% TODO: Grafik einfügen: Rechte-Hand-Regel für a x b.

\subsubsection{Eigenschaften}

\begin{itemize}
     \item $\vec a \times \vec b = \vec 0 \iff \vec a$ und $\vec b$ sind kollinear
     \item $\vec a \times \vec a = \vec 0$
     \item Antikommutativ:
           \[
                \vec a \times \vec b = -(\vec b \times \vec a)
           \]
     \item Distributiv in beiden Argumenten:
           \[
                \vec a \times (\vec b+\vec c)=\vec a \times \vec b + \vec a \times \vec c
           \]
     \item Verträglich mit Skalaren:
           \[
                \lambda(\vec a\times\vec b)=(\lambda\vec a)\times\vec b=\vec a\times(\lambda\vec b)
           \]
     \item Nicht assoziativ:
           \[
                \vec a\times(\vec b\times\vec c)\neq(\vec a\times\vec b)\times\vec c
           \]
\end{itemize}

\subsubsection{Berechnung über Komponenten}
\[
     \vec a \times \vec b =
     \begin{pmatrix}
          a_1 \\
          a_2 \\
          a_3
     \end{pmatrix} \times
     \begin{pmatrix}
          b_1 \\
          b_2 \\
          b_3
     \end{pmatrix}
     =
     \begin{pmatrix}
          a_2b_3-a_3b_2 \\
          a_3b_1-a_1b_3 \\
          a_1b_2-a_2b_1
     \end{pmatrix}
\]

\subsubsection{Geometrische Bedeutung}

Der Betrag des Vektorproduktes entspricht dem Flächeninhalt des von $\vec a$ und $\vec b$ aufgespannten Parallelogramms:

\[
     A = |\vec a \times \vec b|
\]

% TODO: Grafik einfügen: Parallelogramm mit Seiten a,b und Fläche |a x b|.